\clearpage\fancyhead[L]{\small JUNO.}\hypertarget{native-335-lower}{}\label{native-335-lower}
\begin{center}\textbf{JUNO.}\par\vspace{4pt}\hrule\vspace{4pt}\fontsize{10}{12}\selectfont Monatliche Correspondenz. Vol. XXIV. pp. 186\,..\,188. August 1811.\par\vspace{4pt}\hrule\vspace{4pt}\end{center}

\par\indent In the September 1808 issue of this journal [p. 307 of this volume] we gave our readers the elements of \emph{Juno}, which Professor \textsc{Gauss} had revised for the eighth time. Since then another revision had hardly been practicable, as the fifth opposition of \emph{Juno}, which occurred in January 1810, was observed nowhere. It was therefore all the more important not to miss this year’s opposition in April. The planet was so unusually faint that observing it was especially difficult, and Professor \textsc{Gauss} is probably the only one who succeeded in observing this opposition. Here the planet was never visible, even with the high light-gathering power of the Seeberg observatory’s meridian telescope. This makes the few observations that Professor \textsc{Gauss} was able to make with the ring micrometer alone all the more valuable, given the planet’s extreme faintness. For comparison the stars used were $\varphi$ Librae and several other nearby stars from the \emph{Histoire céleste}. Once their positions are determined more precisely, the following results will still need a small correction.
\Needspace{10\baselineskip}\begin{center}\fontsize{10}{12}\selectfont\setlength{\tabcolsep}{4pt}\renewcommand{\arraystretch}{1.15}
\begin{tabular}{l|c|c|c}
1811 & \shortstack{Mean time\\in Göttingen} & \shortstack{Right\\ascension} & \shortstack{Southern\\declination} \\
\hline
April 22 & $9^{\mathrm u}\,51^{\mathrm m}\,35^{\mathrm s}$ & $216^\circ\,41'\,50''1$ & $0^\circ\,58'\,16''$ \\
\phantom{April }24 & $10\;\;32\;\;55$ & $216\;\;17\;\;58.7$ & $0\;\;46\;\;50$ \\
\phantom{April }25 & $10\;\;19\;\;22$ & $216\;\;6\;\;3.6$ & $0\;\;40\;\;4$ \\
\phantom{April }26 & $10\;\;17\;\;30$ & $215\;\;54\;\;7.3$ & $0\;\;33\;\;51$ \\
\end{tabular}
\end{center}

\SourceLeaf{336}\Needspace{12\baselineskip}\par\indent The declination on 22 April is doubtful, and the one on 24 April is also less reliable than the two that follow, when the planet was in a more favorable position. From these observations Professor \textsc{Gauss} found the following result for the opposition.
\begin{center}\fontsize{10}{12}\selectfont\setlength{\tabcolsep}{4pt}\renewcommand{\arraystretch}{1.15}
\begin{tabular}{l}
$\text{\AstroSymbols ☍}$ 1811. 24 April $19^{\mathrm u}\,20^{\mathrm m}\,12^{\mathrm s}$ Mean time in Göttingen \\
$214^\circ\,8'\,48''3$ true heliocentric longitude \\
$12\;\;55\;\;2.0$ true geocentric latitude, north \\
\end{tabular}
\end{center}


\Needspace{20\baselineskip}\par\indent The large discrepancy between this position and the elements based on the first four oppositions now reveals in Juno too the effect of the perturbations caused particularly by \emph{Jupiter}. Calculating them for Juno will be as enormous a task as for Pallas. It is only fair that the work should first be finished for Pallas, which was discovered earlier. For the time being, therefore, Professor \textsc{Gauss} has been content to base new elliptic elements on the oppositions of 1806, 1807, 1808, 1811. We give them below.

\begin{center}\fontsize{10}{12}\selectfont\setlength{\tabcolsep}{4pt}\renewcommand{\arraystretch}{1.15}\begin{tabular}{llr}
\shortstack[l]{Epoch of mean longitude 1811\\Göttingen meridian} & .\;.\;.\;.\;. & $177^\circ\,48'\,1''8$\\
Mean daily tropical motion &  & $813''2486$\\
Longitude of perihelion 1811 & .\;.\;. & $53^\circ\,14'\,32''4$\\
Longitude of the ascending node 1811 &  & $171^\circ\,9'\,13''5$\\
Inclination of the orbit & .\;.\;.\;.\;.\;.\;. & $13\;4\;27.0$\\
Eccentricity & .\;.\;.\;.\;.\;. & $=\mathrm{Sin}\;14\;44\;9.1$\\
Logarithm of the semimajor axis &  & $=0.4265711.$\\
\end{tabular}\end{center}

\par\indent Using these elements, Mr. \textsc{Wachter}, a skilled pupil of Professor \textsc{Gauss}, calculated the following ephemeris for the motion of \emph{Juno} in 1812. [from 23 Feb. to 5 Nov.]
